% AA-19 | Algebra abstrakcyjna | Iteracja X
\chapter[Ciała i ciała ułamków]{Ciała, ciała ułamków i konstrukcje algebraiczne liczb}
\label{aa:ch:fields-fractions}
\setlength{\parskip}{2pt}

\section{Dlaczego pojawia się pojęcie ciała?}
W rozdziałach o pierścieniach pokazaliśmy, że nie każdą niezerową
liczbę całkowitą można odwrócić wewnątrz $\mathbb Z$. Dodanie
odwrotności mnożeniowych prowadzi do $\mathbb Q$, ale konstrukcja
musi zachować wcześniejsze działania i nie może utożsamić
różnych liczb całkowitych. Jednocześnie ta sama idea działa
dla dowolnej dziedziny całkowitości, na przykład dla $K[t]$.
Właśnie to wyjaśnia konstrukcja ciała ułamków.

\begin{definicja}[Ciało, podciało i homomorfizm ciał]
\label{aa:def:field}
\emph{Ciałem} nazywamy niezerowy pierścień przemienny $K$ z jedynką,
w którym każdy element $a\in K\setminus\{0\}$ ma odwrotność
mnożeniową $a^{-1}\in K$. Zbiór $K^\times=K\setminus\{0\}$
jest wówczas grupą abelową względem mnożenia:
jeśli $a,b\ne0$ i $ab=0$, pomnożenie przez $a^{-1}$
dałoby $b=0$, sprzeczność. Jedynka i odwrotności
należą do $K^\times$, a łączność i przemienność
mnożenia dziedziczymy z $K$.
\emph{Podciało} $F\subseteq K$ jest podzbiorem zawierającym
$0,1$, zamkniętym na dodawanie, odejmowanie, mnożenie
oraz odwrotności elementów niezerowych. Dla ciał $K,L$ \emph{homomorfizmem ciał}
nazywamy odwzorowanie $f:K\to L$ zachowujące
$0,1,+,\cdot$. Zachowuje ono również odwrotności:
dla $a\ne0$ równość $f(a)f(a^{-1})=f(1)=1$
pokazuje, że $f(a^{-1})=f(a)^{-1}$.
\end{definicja}

\begin{stwierdzenie}[Homomorfizm ciał jest różnowartościowy]
\label{aa:st:field-embedding-injective}
Niech $K,L$ będą ciałami, a $f:K\to L$
homomorfizmem ciał. Wówczas $f$ jest różnowartościowy.
\end{stwierdzenie}
\begin{proof}
Niech $a\in\ker f$ i załóżmy, że $a\ne0$. Wówczas
$1_K=aa^{-1}$, a z zachowywania mnożenia mamy
$1_L=f(1_K)=f(a)f(a^{-1})=0$, co przeczy definicji ciała.
Zatem $\ker f=\{0\}$. Jeśli $f(x)=f(y)$, to
$f(x-y)=0$ i stąd $x-y=0$, czyli $x=y$.
\end{proof}

\begin{uwaga}[Nie mylić ciała z pierścieniem z dzieleniem]
W ciele zakładamy przemienność mnożenia. W nieprzemiennym
pierścieniu z dzieleniem każdy element niezerowy również
jest odwracalny, ale może nie komutować z innymi.
Przykładem jest algebra kwaternionów Hamiltona $\mathbb H$,
która jest pierścieniem z dzieleniem, lecz nie jest ciałem
w przyjętym tu, przemiennym znaczeniu.
\end{uwaga}

\section{Ciało ułamków dowolnej dziedziny całkowitości}
\begin{definicja}[Relacja ułamków formalnych]
\label{aa:def:fractions-relation}
Niech $R$ będzie dziedziną całkowitości, a
$S=R\setminus\{0\}$. Na $R\times S$ definiujemy
$(a,b)\sim(c,d)$ wtedy i tylko wtedy, gdy $ad=bc$.
Klasę $(a,b)$ oznaczamy $a/b$; na razie jest to
symbol formalny, a nie dzielenie wykonane w $R$.
\end{definicja}

\begin{lemat}[Poprawność relacji równoważności]
\label{aa:lem:fractions-equivalence}
Niech $R$ będzie dziedziną całkowitości i
$S=R\setminus\{0\}$. Relacja z
Definicji~\ref{aa:def:fractions-relation} jest relacją
równoważności na $R\times S$.
\end{lemat}
\begin{proof}
Zwrotność wynika z $ab=ba$. Jeśli $ad=bc$, to
$cb=da$, więc zachodzi symetria. Dla przechodniości
załóżmy $(a,b)\sim(c,d)$ oraz $(c,d)\sim(e,f)$.
Oznacza to $ad=bc$ oraz $cf=de$. Pierwszą równość
mnożymy przez $f$, a drugą przez $b$:
$adf=bcf=bde$. Stąd $d(af-be)=0$. Ponieważ
$d\in S$, brak dzielników zera w $R$ daje $af=be$.
Jest to dokładnie warunek $(a,b)\sim(e,f)$.
\end{proof}

\begin{twierdzenie}[Istnienie ciała ułamków]
\label{aa:tw:fractions-field}
Niech $R$ będzie dziedziną całkowitości,
$S=R\setminus\{0\}$ i niech $\sim$ będzie relacją
z Definicji~\ref{aa:def:fractions-relation}. Zbiór klas
$K=(R\times S)/{\sim}$ z działaniami
\[
 \frac ab+\frac cd=\frac{ad+bc}{bd},\qquad
 \frac ab\cdot\frac cd=\frac{ac}{bd}
\]
jest ciałem. Zero i jedynkę reprezentują odpowiednio
$0/1$ i $1/1$, a dla $a/b\ne0/1$ odwrotnością jest $b/a$.
Odwzorowanie $\iota:R\to K$, $a\mapsto a/1$,
jest różnowartościowym homomorfizmem pierścieni.
\end{twierdzenie}
\begin{proof}
\emph{(1) Zamkniętość.} Ponieważ $b,d$ są niezerowe,
a $R$ jest dziedziną całkowitości, również $bd\ne0$.
Oba prawe wyrażenia wyznaczają więc poprawne pary.

\emph{(2) Niezależność od reprezentantów.}
Niech $a/b=a'/b'$ oraz $c/d=c'/d'$,
czyli $ab'=a'b$ i $cd'=c'd$. Wówczas
$ac\,b'd'=(ab')(cd')=(a'b)(c'd)=a'c'\,bd$,
co dowodzi poprawności mnożenia.
W przypadku dodawania sprawdzamy równość
$(ad+bc)b'd'=(a'd'+b'c')bd$:
po rozwinięciu pierwszego składnika mamy
$ad b'd'=ab'\,dd'=a'b\,dd'=a'd'\,bd$;
dla drugiego składnika analogicznie
$bc b'd'=bb'\,cd'=bb'\,c'd=b'c'\,bd$.
Dodając obie równości dostajemy żądany wynik.

\emph{(3) Aksjomaty pierścienia.} Przemienność i łączność
mnożenia wynikają z odpowiednich własności $R$:
$(a/b)(c/d)(e/f)=ace/(bdf)$, niezależnie od nawiasowania.
Dodawanie jest przemienne, a dla trzech ułamków
obydwa nawiasowania sumy $a/b$, $c/d$, $e/f$ dają
$(adf+cbf+ebd)/(bdf)$. Każdy składnik licznika
powstaje przez pomnożenie odpowiedniego licznika
przez dwa pozostałe mianowniki. Element $0/1$ jest zerem, gdyż
$(a/b)+(0/1)=a/b$. Przeciwieństwem $a/b$
jest $(-a)/b$, ponieważ ich suma jest $0/b^2=0/1$.
Ponadto $(a/b)(1/1)=a/b$, więc $1/1$ jest jedynką.
Różnią się one, gdyż $0/1=1/1$ oznaczałoby $0=1$
w dziedzinie $R$.
Rozdzielność sprawdzamy bez skrótów:
$(a/b)((c/d)+(e/f))=a(cf+ed)/(bdf)$
i $(a/b)(c/d)+(a/b)(e/f)=(acf+aed)/(bdf)$.
Licznikami są w obu przypadkach $acf+aed$.

\emph{(4) Odwrotności.} Warunek $a/b=0/1$ oznacza $a=0$.
Jeśli więc $a/b\ne0/1$, to $a\ne0$ i $b/a$ jest
dopuszczalnym ułamkiem. Iloczyn wynosi
$(a/b)(b/a)=ab/(ab)=1/1$, ponieważ $ab\ne0$
i $(ab)\cdot1=1\cdot(ab)$.
Ponadto wzór nie zależy od reprezentanta: jeśli
$a/b=c/d\ne0/1$, to $a,c\ne0$, a równość
$ad=bc$ jest także warunkiem równoważności par
$(b,a)$ i $(d,c)$. Zatem $b/a=d/c$.

\emph{(5) Zanurzenie.} Wzory na działania dają
$\iota(a+c)=\iota(a)+\iota(c)$ i
$\iota(ac)=\iota(a)\iota(c)$; ponadto
$\iota(1)=1/1$. Jeśli $\iota(a)=\iota(c)$,
to $a\cdot1=c\cdot1$, zatem $a=c$.
\end{proof}


\begin{definicja}[Ciało ułamków]
\label{aa:def:field-of-fractions}
Niech $R$ będzie dziedziną całkowitości.
Ciało skonstruowane dla $R$ w
Twierdzeniu~\ref{aa:tw:fractions-field} nazywamy
\emph{ciałem ułamków} dziedziny $R$ i oznaczamy
$\operatorname{Frac}(R)$.
\end{definicja}

\begin{twierdzenie}[Własność uniwersalna ciała ułamków]
\label{aa:tw:fractions-universal}
Niech $R$ będzie dziedziną całkowitości, $L$ ciałem,
a $f:R\to L$ różnowartościowym homomorfizmem pierścieni.
Istnieje dokładnie jeden homomorfizm ciał
$\overline f:\operatorname{Frac}(R)\to L$ taki, że
$\overline f(a/1)=f(a)$ dla każdego $a\in R$.
Jest nim odwzorowanie $\overline f(a/b)=f(a)f(b)^{-1}$.
\end{twierdzenie}
\begin{proof}
Ponieważ $b\ne0$ oraz $f$ jest różnowartościowy,
$f(b)\ne0$ i $f(b)^{-1}$ istnieje. Jeśli $a/b=c/d$,
to $ad=bc$, a więc $f(a)f(d)=f(b)f(c)$.
Mnożąc przez $f(b)^{-1}f(d)^{-1}$ dostajemy
$f(a)f(b)^{-1}=f(c)f(d)^{-1}$. Wzór nie zależy
od reprezentantów. Dla $a/b,c/d$ otrzymujemy
\[
\overline f\left(\frac ab+\frac cd\right)
=\frac{f(a)f(d)+f(b)f(c)}{f(b)f(d)}
=\frac{f(a)}{f(b)}+\frac{f(c)}{f(d)},
\]
a dla mnożenia, korzystając z przemienności w $L$,
mamy
\[
\overline f\left(\frac ab\cdot\frac cd\right)
=f(ac)f(bd)^{-1}
=\bigl(f(a)f(b)^{-1}\bigr)
 \bigl(f(c)f(d)^{-1}\bigr).
\]
Ponadto $\overline f(0/1)=f(0)=0$ i
$\overline f(1/1)=f(1)=1$, więc jest to homomorfizm ciał.
Dla jednoznaczności zauważmy, że $b/1$ jest odwracalne
w $\operatorname{Frac}(R)$ i że
$a/b=(a/1)(b/1)^{-1}$. Każdy homomorfizm ciał
przedłużający $f$ musi zachować ten iloczyn oraz odwrotność,
a zatem musi działać według podanego wzoru.
Różnowartościowość $f$ jest niezbędna dla tego wzoru:
bez niej jakiś niezerowy mianownik $b$ mógłby mieć
$f(b)=0$, którego nie można odwrócić w $L$.
Na przykład redukcja $\mathbb Z\to\mathbb F_p$
nie przedłuża się do homomorfizmu ciał
$\mathbb Q\to\mathbb F_p$: odwrotność liczby
$p/1$ musiałaby być wysłana na odwrotność zera.
\end{proof}

\begin{przyklad}[Liczby wymierne jako ciało ułamków]
\label{aa:prz:construct-rationals}
Wybierzmy $R=\mathbb Z$. Wówczas otrzymujemy
\[
\mathbb Q\cong\operatorname{Frac}(\mathbb Z)
=(\mathbb Z\times(\mathbb Z\setminus\{0\}))/\!\sim,
\qquad (a,b)\sim(c,d)\iff ad=bc.
\]
Na przykład $(2,3)\sim(4,6)$, lecz $(2,3)\not\sim(3,4)$.
Równość $1/2+1/3=5/6$ wynika z formalnego wzoru
$(1\cdot3+1\cdot2)/(2\cdot3)$, a nie z wcześniej
założonej arytmetyki ułamków. Różne pary reprezentują tę
samą liczbę dokładnie wtedy, gdy zachodzi równość
iloczynów krzyżowych. Element $-5/7$ można zapisać
jako $5/(-7)$, ponieważ $(-5)\cdot(-7)=5\cdot7$.
W szczególności wcześniejsza konstrukcja $\mathbb Z$
z $\mathbb N_0\times\mathbb N_0$ oraz obecna konstrukcja
$\mathbb Q$ tworzą łańcuch zanurzeń
$\mathbb N_0\hookrightarrow\mathbb Z\hookrightarrow\mathbb Q$.
\end{przyklad}

\begin{przyklad}[Funkcje wymierne]
\label{aa:prz:rational-functions}
Dla dowolnego ciała $K$ pierścień $K[t]$ jest dziedziną
całkowitości, a jego ciało ułamków oznaczamy $K(t)$.
Elementem $K(t)$ jest klasa pary wielomianów $(p,q)$,
$q\ne0$, którą zapisujemy $p(t)/q(t)$.
Przykładowo $t/(t+1)$ i $t^2/(t^2+t)$ są równe
w $K(t)$, gdyż $t(t^2+t)=t^2(t+1)$;
oba mianowniki są niezerowymi wielomianami.
Wspólna klasa tych ułamków nie należy do podpierścienia
$K[t]\subseteq K(t)$. Gdyby istniał wielomian $h\in K[t]$
taki, że $t/(t+1)=h$, to w $K[t]$ mielibyśmy
$t=h(t)(t+1)$; podstawienie $t=-1$ dałoby $-1=0$,
co przeczy $1\ne0$ w ciele $K$.
Jeżeli $p/q=0$, to $p=0$, gdyż $K[t]$ nie ma
dzielników zera.
\end{przyklad}

\begin{przyklad}[Dlaczego nie wystarczy pierścień z dzielnikami zera]
\label{aa:prz:frac-zero-divisors}
W $R=\mathbb Z/6\mathbb Z$ mamy $[2][3]=0$ mimo
$[2]\ne0$ i $[3]\ne0$. Gdyby można było zanurzyć
$R$ w ciało $K$, to obrazy tych dwóch elementów byłyby
niezerowe, lecz ich iloczyn równy zeru. Jest to niemożliwe
w ciele. Nie istnieje więc ciało zawierające $R$ jako
podpierścień. Można badać \emph{całkowity pierścień ułamków}
po odwróceniu wszystkich elementów niebędących dzielnikami
zera, lecz nie jest on na ogół ciałem.
\end{przyklad}

\section{Charakterystyka i ciało pierwsze}
\begin{definicja}[Charakterystyka pierścienia]
\label{aa:def:field-characteristic}
Dla pierścienia $R$ z jedynką i $1_R\ne0_R$
\emph{charakterystyką} nazywamy najmniejszą dodatnią
liczbę całkowitą $n$, dla której $n1_R=0_R$, jeśli
taka istnieje; w przeciwnym przypadku charakterystyka
wynosi zero. Zapis $n1_R$ oznacza sumę $n$ kopii $1_R$.
\end{definicja}

\begin{twierdzenie}[Charakterystyka ciała]
\label{aa:tw:field-characteristic-prime}
Charakterystyka dowolnego ciała $K$ jest równa zeru
albo liczbie pierwszej $p$.
\end{twierdzenie}
\begin{proof}
Niech $n>0$ będzie charakterystyką $K$.
Gdyby $n=ab$ dla $1<a,b<n$, to z minimalności $n$
wynikałoby $a1_K\ne0$ i $b1_K\ne0$.
Jednak $(a1_K)(b1_K)=(ab)1_K=n1_K=0$.
Jest to niemożliwe, ponieważ ciało nie ma dzielników zera.
Zatem $n$ nie jest złożone, a minimalność i $1_K\ne0$
wykluczają $n=1$. Stąd $n$ jest liczbą pierwszą.
\end{proof}

\begin{twierdzenie}[Najmniejsze podciało]
\label{aa:tw:prime-subfield}
Każde ciało $K$ zawiera dokładnie jedno najmniejsze podciało.
Jest ono izomorficzne z $\mathbb Q$, jeśli
$\operatorname{char}K=0$, oraz z $\mathbb F_p=\mathbb Z/p\mathbb Z$,
jeśli $\operatorname{char}K=p>0$.
\end{twierdzenie}
\begin{proof}
Rozważmy homomorfizm pierścieni
$\phi:\mathbb Z\to K$, $n\mapsto n1_K$.
Jego jądro jest ideałem $\mathbb Z$, a więc ma
postać $m\mathbb Z$ dla pewnego $m\ge0$.
Jeśli charakterystyka jest zerowa, to $m=0$ i $\phi$
jest różnowartościowy. Na mocy
Twierdzenia~\ref{aa:tw:fractions-universal} istnieje
jednoznaczne zanurzenie $\mathbb Q\to K$.
Jeżeli charakterystyka jest dodatnia, to z poprzedniego
twierdzenia wynika $m=p$; pierwsze twierdzenie o
izomorfizmie daje zanurzenie
$\mathbb Z/p\mathbb Z\to K$. Iloraz ten jest ciałem,
ponieważ dla każdego $a$ niepodzielnego przez $p$
tożsamość Bézouta $ua+vp=1$ daje odwrotność klasy $[u]$.

Każde podciało $F\subseteq K$ zawiera $1_K$, a zatem
wszystkie wielokrotności $n1_K$ i odwrotności wszystkich
niezerowych spośród nich. W charakterystyce zero zawiera
też wszystkie ilorazy $(m1_K)(n1_K)^{-1}$ dla $n\ne0$,
czyli cały obraz $\mathbb Q$. Musi więc zawierać opisany obraz
$\mathbb Q$ lub $\mathbb F_p$. Obraz sam jest podciałem,
stąd jest najmniejszy i jedyny.
\end{proof}

\begin{twierdzenie}[Endomorfizm Frobeniusa]
\label{aa:tw:field-frobenius}
Niech $K$ będzie ciałem charakterystyki $p>0$.
Odwzorowanie $F:K\to K$, $F(x)=x^p$, jest
różnowartościowym homomorfizmem ciał. Jeśli $K$
jest skończone, $F$ jest automorfizmem.
\end{twierdzenie}
\begin{proof}
Wzór dwumianowy daje
$(x+y)^p=\sum_{j=0}^p\binom pj x^j y^{p-j}$.
Dla $1\le j\le p-1$ współczynnik $\binom pj$
jest podzielny przez $p$: w równości całkowitoliczbowej
$\binom pj\,j!(p-j)!=p!$ żaden z czynników $j!$ ani $(p-j)!$ nie jest
podzielny przez $p$.
Ponieważ $p$ jest liczbą pierwszą, z lematu Euklidesa wynika
$p\mid\binom pj$.
W charakterystyce $p$ pośrednie składniki znikają,
więc $(x+y)^p=x^p+y^p$. Równość
$(xy)^p=x^py^p$ jest oczywista z przemienności.
Ponieważ również $F(1)=1^p=1$,
$F$ jest homomorfizmem ciał, a
Twierdzenie~\ref{aa:st:field-embedding-injective}
daje jego różnowartościowość.
Na zbiorze skończonym każda funkcja różnowartościowa
jest surjektywna, więc $F$ jest automorfizmem.
\end{proof}

\begin{przyklad}[Frobenius nie zawsze jest na]
\label{aa:prz:frobenius-not-surjective}
Niech $p$ będzie liczbą pierwszą. W ciele
$K=\mathbb F_p(t)$ element $t$ nie jest
$p$-tą potęgą żadnego elementu. Gdyby
$(a(t)/b(t))^p=t$ dla wielomianów $a,b\in\mathbb F_p[t]$,
$b\ne0$, to $a(t)\ne0$ i $a(t)^p=t b(t)^p$.
Po porównaniu stopni niezerowych wielomianów
otrzymalibyśmy $p\deg a=1+p\deg b$, co jest
niemożliwe modulo $p$. Endomorfizm Frobeniusa
nie jest zatem surjektywny w tym nieskończonym ciele.
\end{przyklad}

\section{Dwie algebraiczne konstrukcje ciała liczb zespolonych}
W tej sekcji \emph{przyjmujemy ciało liczb rzeczywistych}
$\mathbb R$ wraz ze znanym porządkiem i jego własnościami.
Nie konstruujemy tu $\mathbb R$: konstrukcja ciała
uporządkowanego zupełnego należy do skryptu analizy.
Pokazujemy natomiast, że operacje na parach liczb
rzeczywistych i odpowiedni pierścień ilorazowy prowadzą
do tego samego ciała.

\begin{definicja}[Liczby zespolone jako pary]
\label{aa:def:complex-as-pairs}
Na zbiorze $\mathbb R^2$ określamy
\[
(a,b)+(c,d)=(a+c,b+d),\qquad
(a,b)(c,d)=(ac-bd,ad+bc).
\]
Element neutralny dodawania to $(0,0)$, a jedynka
mnożenia to $(1,0)$. Oznaczamy $i=(0,1)$.
\end{definicja}

\begin{twierdzenie}[Struktura ciała na $\mathbb R^2$]
\label{aa:tw:complex-pairs-field}
Działania z Definicji~\ref{aa:def:complex-as-pairs}
czynią $\mathbb R^2$ ciałem,
w którym $i^2=(-1,0)$ oraz
$(a,b)=(a,0)+(b,0)i$.
Dla $(a,b)\ne(0,0)$ odwrotność jest równa
\[
(a,b)^{-1}=\left(\frac a{a^2+b^2},
 -\frac b{a^2+b^2}\right).
\]
\end{twierdzenie}
\begin{proof}
Wzory na dodawanie dają grupę abelową z przeciwnym
elementem $(-a,-b)$. Przemienność mnożenia
wynika z $ac=ca$, $bd=db$ i $ad+bc=cb+da$.
Dla łączności mnożenia obliczmy
$((a,b)(c,d))(e,f)$: jego pierwsza współrzędna to
$(ac-bd)e-(ad+bc)f=ace-bde-adf-bcf$,
a druga to $(ac-bd)f+(ad+bc)e=acf-bdf+ade+bce$.
Obliczając $(c,d)(e,f)=(ce-df,cf+de)$, otrzymujemy
z $(a,b)((c,d)(e,f))$ pierwszą współrzędną
$a(ce-df)-b(cf+de)=ace-adf-bcf-bde$
i drugą $a(cf+de)+b(ce-df)=acf+ade+bce-bdf$.
Są to te same wyrażenia, więc mnożenie jest łączne.
Dla rozdzielności po lewej stronie porównajmy
$(a,b)((c,d)+(e,f))$ z sumą dwóch iloczynów:
w obu przypadkach otrzymujemy parę
$(ac+ae-bd-bf,ad+af+bc+be)$.
Rozdzielność po prawej wynika następnie z przemienności. Bezpośredni rachunek daje
$(1,0)(a,b)=(a,b)$, $(0,1)^2=(-1,0)$
oraz $(a,0)+(b,0)(0,1)=(a,b)$.

Jeżeli $(a,b)\ne(0,0)$, to $a^2+b^2>0$, a więc
mianownik w podanym wzorze nie jest zerem. Iloczyn
$(a,b)(a,-b)=(a^2+b^2,0)$; podzielenie współrzędnych
przez $a^2+b^2$ daje wskazaną odwrotność.
Ponieważ istnieją odwrotności wszystkich elementów
niezerowych, otrzymaliśmy ciało.
\end{proof}

\begin{twierdzenie}[Liczby zespolone jako pierścień ilorazowy]
\label{aa:tw:complex-quotient}
Wielomian $X^2+1\in\mathbb R[X]$ jest nierozkładalny.
Pierścień $\mathbb R[X]/(X^2+1)$ jest ciałem i istnieje
izomorfizm ciał
\[
\mathbb R[X]/(X^2+1)\longrightarrow\mathbb R^2,
\qquad [a+bX]\longmapsto(a,b).
\]
Izomorfizm jest wyznaczony przez przesłanie klasy
$[X]$ na element $i=(0,1)$, dla którego $i^2=(-1,0)$.
\end{twierdzenie}
\begin{proof}
Wielomian drugiego stopnia nad ciałem jest rozkładalny
wtedy i tylko wtedy, gdy ma pierwiastek w tym ciele:
rozkład nietrywialny ma czynniki liniowe, a pierwiastek
$x$ daje rozkład $f(X)=(X-x)g(X)$ na mocy
twierdzenia o reszcie.
Gdyby $X^2+1$ miał pierwiastek $x\in\mathbb R$,
to $x^2=-1$, co przeczyłoby $x^2\ge0$.
Zatem $X^2+1$ jest nierozkładalny. Ponieważ $\mathbb R[X]$ jest dziedziną euklidesową,
dla każdego $g\notin(X^2+1)$ największy wspólny dzielnik
$g$ i $X^2+1$ jest jednostką: nierozkładalność $X^2+1$
wyklucza nieodwracalny wspólny dzielnik.
Tożsamość Bézouta daje wówczas wielomiany $u,v$ z
$ug+v(X^2+1)=1$, czyli $[u][g]=[1]$ w ilorazie.
Każda niezerowa klasa jest odwracalna, zatem iloraz
jest ciałem, a ideał $(X^2+1)$ jest maksymalny.

Dzielenie z resztą przez unormowany wielomian
$X^2+1$ pokazuje, że każda klasa ma reprezentanta
postaci $a+bX$. Jest on jednoznaczny: różnica dwóch
takich reprezentantów o stopniu mniejszym niż $2$
nie może być niezerową wielokrotnością $X^2+1$. W ilorazie
$[X]^2=[-1]$, więc
$(a+b[X])(c+d[X])=(ac-bd)+(ad+bc)[X]$.
Wskazany wzór zachowuje działania i jedynkę.
Jest różnowartościowy dzięki jednoznaczności reszty,
a surjektywny, bo każda para $(a,b)$ jest obrazem
klasy $a+b[X]$. Otrzymaliśmy izomorfizm ciał. Powstałe ciało
$\mathbb R^2$ z powyższymi działaniami oznaczamy
$\mathbb C$.
\end{proof}

\begin{przyklad}[Zespolone równanie kwadratowe bez analizy]
\label{aa:prz:complex-square-root}
W skonstruowanym ciele $\mathbb C$ równanie
$z^2+1=0$ ma dwa rozwiązania: $i$ i $-i$.
Pierwszy wniosek to $i^2=-1$. Ponieważ
$z^2+1=(z-i)(z+i)$ i $\mathbb C$ nie ma dzielników
zera, każde rozwiązanie musi być równe $i$ lub $-i$.
Dwa rozwiązania są różne, gdyż ciało uporządkowane
$\mathbb R$ ma charakterystykę zero:
$n\cdot1>0$ dla każdego $n\ge1$, zatem $2i\ne0$.
\end{przyklad}

\section{Granice opisu konstrukcyjnego}
Ta sama ogólna idea pojawi się w kolejnym rozdziale:
iloraz $K[X]/(f)$ pozwala \emph{dołączyć pierwiastek}
nierozkładalnego wielomianu $f$. Iloraz przez ideał generowany przez wielomian rozkładalny
nie musi być ciałem.
Dla przykładu $K[X]/(X^2)$ zawiera niezerową klasę
$[X]$ o kwadracie równym zeru. Nie jest więc ciałem.
Taka konfrontacja przykładów z twierdzeniem wyjaśnia,
dlaczego nierozkładalność wielomianu jest istotnym
założeniem konstrukcji rozszerzeń ciał.

\setlength{\parskip}{4pt}
